\chapter{Metodología de evaluación experimental}
\label{chap:metodologia-evaluacion}

Este capítulo explica cómo se comparan M5, S1 y H y qué evidencia se recoge para
responder a las preguntas de investigación. La evaluación se limita a gestiones
administrativas simuladas con perfiles y agendas sintéticos. No toma decisiones
clínicas ni representa una prueba en un sistema sanitario real.

El estudio compara las tres arquitecturas con dos modelos:
\texttt{openai/gpt-oss-20b} servido por Groq \citep{openai2025gptoss} y
\texttt{gemini-3.5-flash-lite} servido por Google AI Studio
\citep{gemini35FlashLiteDocs}. La comparación principal se hace entre
arquitecturas con el mismo modelo. La comparación entre modelos es secundaria y
no separa el efecto del modelo del de su proveedor.

\section{Preguntas de investigación y comparación}
\label{sec:metodo-preguntas}

\subsection{Fundamento del procedimiento}

La evaluación distingue el resultado de la gestión, los recursos empleados y
la calidad de la conversación. Consultar el estado final de la agenda permite
comprobar una reserva, pero no explica si la persona tuvo que repetir la
información o si el asistente comunicó bien los pasos. Por eso se combinan
comprobaciones del estado y de los eventos con métricas de consumo y revisión
de las conversaciones.

La Tabla~\ref{tab:fuentes-metodo} relaciona estas decisiones con la literatura.
Los trabajos citados orientan el procedimiento. Los casos, el indicador de
éxito integral y la rúbrica se definen para las gestiones de CitaLab.

\begin{longtable}{L{0.25\textwidth}L{0.38\textwidth}L{0.26\textwidth}}
\caption{Fuentes que fundamentan las decisiones de evaluación.}
\label{tab:fuentes-metodo}\\
\toprule
\textbf{Fuente} & \textbf{Decisión adoptada} & \textbf{Límites} \\
\midrule
\endfirsthead
\toprule
\textbf{Fuente} & \textbf{Decisión adoptada} & \textbf{Límites} \\
\midrule
\endhead
\(\tau\)-bench y \(\tau^2\)-bench \citep{yao2024taubench,barres2025tau2}
& Evaluar diálogo, operaciones, estado final y restricciones con un usuario simulado que solo ve información pública.
& Son otros dominios. El usuario simulado requiere validación propia. \\
AgentBoard y MAST \citep{ma2024agentboard,cemri2025mast}
& Localizar la primera desviación y su propagación en las trazas.
& El orden de eventos no identifica causalidad. \\
LaMP y diálogo multivuelta \citep{salemi2024lamp,laban2025multiturn}
& Separar uso del perfil y calidad de la respuesta. Probar información progresiva y correcciones.
& No se infiere memoria longitudinal ni satisfacción de pacientes. \\
PARADISE \citep{walker1997paradise}
& Informar cumplimiento de la tarea, turnos y esfuerzo por separado.
& No se adopta su modelo de satisfacción. \\
AI Agents That Matter \citep{kapoor2024agentsthatmatter}
& Comparar utilidad y coste, congelar versiones y declarar la exposición del corpus.
& Corpus final nuevo. Categorías y requisitos conocidos durante el desarrollo. \\
Fiabilidad de agentes \citep{rabanser2026reliability}
& Distinguir consistencia, robustez y seguridad de una tasa media.
& Pocas repeticiones describen estabilidad observada, no fiabilidad general. \\
Significación estadística \citep{dror2018significance}
& Respetar el emparejamiento en las comparaciones.
& Pocas unidades independientes. \\
Jueces y jurados LLM \citep{zheng2023llmjudge,verga2024poll,panickssery2024selfpreference,bavaresco2025judgebench}
& Puntuar todas las conversaciones con dos jueces de proveedores distintos y contrastarlos con una revisión humana ciega (apartado~\ref{sec:metodo-mixta}).
& Su validez depende de la tarea. Sin revisores externos. \\
Inferencia potenciada por predicción \citep{angelopoulos2023ppi,angelopoulos2023ppiplus,boyeau2025autoeval}
& Se previó combinar el jurado en todos los episodios con la muestra humana. Se descartó porque la muestra no representa toda la campaña (apartado~\ref{sec:metodo-mixta}).
& Requiere que la muestra represente los episodios a los que se extiende. \\
\bottomrule
\end{longtable}

\subsection{Contraste arquitectónico}

Las preguntas PI1--PI4 de la introducción se concretan así:

\begin{description}[style=nextline]
  \item[PI1 (agentes especializados frente a agente único).] Diferencia pareada
  de éxito integral M5--S1, con su coste en llamadas, tokens y tiempo de inferencia. H se
  añade como referencia de flujo programado.
  \item[PI2 (preferencias).] Adherencia a la preferencia guardada cuando la
  petición deja libertad (B02, R02 y sus contrafactuales) y respeto de la
  petición actual cuando la contradice (B04, M04).
  \item[PI3 (controles).] Violaciones de seguridad (escrituras sin confirmación,
  efectos sobre terceros, cambios colaterales y afirmaciones falsas) y bloqueos
  de gestiones legítimas.
  \item[PI4 (evaluación conjunta).] Éxito integral más calidad de la
  conversación (jurado LLM contrastado con una muestra humana ciega), informados
  por separado.
\end{description}

Dentro de cada modelo, las tres variantes reciben el mismo caso, perfil, reloj,
agenda, permisos y herramientas. Los \emph{prompts} cambian porque forman parte
de cada arquitectura.

\section{Corpus de escenarios}
\label{sec:metodo-corpus}

El corpus contiene veinticinco familias sintéticas, cinco por gestión: reserva
(B01--B05), consulta (Q01--Q05), reprogramación (R01--R05), cancelación
(C01--C05) y recordatorio (M01--M05). En el corpus final los identificadores
llevan el prefijo \texttt{N-} (por ejemplo, \texttt{N-B01}). En este capítulo
se abrevian sin el prefijo. Cada familia define:

\begin{itemize}
  \item el primer mensaje, el idioma (español o inglés) y un estilo de
        comunicación (directo, progresivo, que pide orientación, detallista o
        cambiante).
  \item el objetivo de la persona y un guion condicional (por ejemplo,
        \enquote{cuando veas la primera propuesta, corrige la fecha}).
  \item la agenda inicial, el perfil y, si procede, un evento externo (otra
        persona ocupa el hueco propuesto antes de la confirmación).
  \item el efecto esperado en forma verificable (cita, hueco, canal o
        antelación) y comprobaciones de la respuesta (idioma, datos que deben o
        no deben aparecer).
\end{itemize}

Las familias cubren peticiones claras, información progresiva, preferencias
blandas frente a restricciones expresas, preguntas clínicas incidentales, datos
de otra persona, modalidades no disponibles, correcciones antes de confirmar,
desistimientos y conflictos de concurrencia. Todos los datos son sintéticos.
La Tabla~\ref{tab:diseno-campana} resume cuántas condiciones y episodios aporta
cada bloque de la campaña.

\begin{table}[!htbp]
\centering
\caption{Diseño de la campaña (episodios = unidades \(\times\) 3 arquitecturas \(\times\) 2 modelos).}
\label{tab:diseno-campana}
\begin{tabular}{lrr}
\toprule
Bloque & Unidades & Episodios \\
\midrule
Primera ejecución & 25 familias & 150 \\
Repetición (B01, Q01, R01, C01, M01) & 5 & 30 \\
Contrastes de preferencias (B02, Q02, R02, C02, M02) & 5 & 30 \\
Contrastes de persona (B01, B02, Q01, Q02, R01, M01) & 6 & 36 \\
Casos nuevos (suplemento prospectivo) & 5 & 30 \\
\midrule
Total & 46 & 276 \\
\bottomrule
\end{tabular}
\end{table}

Los contrastes de preferencias cambian un único atributo del perfil (franja,
nivel de detalle o canal) manteniendo el mensaje y la agenda. Permiten comprobar
si la respuesta cambia de forma coherente con el perfil. Los contrastes de
persona mantienen el caso, el perfil y el primer mensaje y cambian una única
dimensión del usuario simulado: familiaridad con el procedimiento (habitual o
primera vez), conocimiento de sus citas, iniciativa al aportar datos,
preferencia por explicaciones paso a paso, familiaridad con los recordatorios y
rechazo público al uso de su perfil. Estudian si las arquitecturas responden
igual de bien a usuarios distintos ante la misma gestión. En los resultados,
los contrastes de preferencias se identifican con \texttt{CF} (por ejemplo,
\texttt{N-CF-R02}), los contrastes de persona con \texttt{UC}
(\texttt{N-UC-B02}) y los casos nuevos con \texttt{SUP} seguido de la letra de
la gestión (\texttt{N-SUP-Q}).

Las conversaciones utilizadas para ajustar las instrucciones pertenecen al
corpus de desarrollo. La campaña final utiliza otro conjunto, con mensajes,
objetivos concretos y agenda propios. Se mantiene el diseño de 276 episodios y
las mismas categorías de gestión. Esta separación evita evaluar únicamente
las conversaciones utilizadas durante el ajuste \citep{kapoor2024agentsthatmatter}.

Los casos finales se redactaron antes de observar sus respuestas. No se
seleccionan por los resultados obtenidos ni se utilizan para modificar las
instrucciones dentro de la campaña. Como corresponde a un corpus construido
para evaluar las gestiones de CitaLab, comparten los requisitos, las
capacidades y las formas de interacción consideradas durante el diseño. Esa
cercanía puede favorecer en cierta medida los resultados, y hay que tenerla en
cuenta al extenderlos a otros contextos.

\section{Usuario simulado}
\label{sec:metodo-simulador}

Cada episodio es una conversación entre el asistente y un usuario simulado,
siguiendo el enfoque de simuladores guiados por un objetivo e información
pendiente \citep{schatzmann2007agenda} con un LLM como generador de los
mensajes \citep{ni2026simulation,mehri2026goalalignment}. El simulador recibe el
idioma, la persona funcional (directa, progresiva, que pide orientación,
detallista o cambiante), el objetivo y un guion condicional. Ve la conversación
y, si hay una propuesta pendiente, la tarjeta que CitaLab mostraría con su botón
Confirmar. No conoce la arquitectura, el estado interno ni el oráculo.

En cada turno responde con una acción: escribir un mensaje, pulsar Confirmar
(solo si hay una propuesta en pantalla) o terminar. Si intenta confirmar sin
propuesta visible, el arnés convierte la acción en un mensaje de texto, como
ocurriría en la interfaz, y lo registra. El episodio termina cuando el simulador
da el objetivo por cumplido o imposible, ante una escalada por emergencia o tras
diez turnos. Un fallo agotado de infraestructura interrumpe el intento como
incidente repetible y conserva la conversación parcial.

\subsection{Elección del modelo simulador}

El usuario simulado utiliza \path{qwen/qwen3.8-27b}, servido por Groq. Es una
familia distinta de GPT-OSS y Gemini, que son los modelos del asistente. Con
ello se intenta reducir la influencia de patrones compartidos entre quien
produce la respuesta y quien continúa el diálogo. Esta elección no elimina
los posibles sesgos del simulador \citep{zheng2023llmjudge}.

La salida estructurada permite indicar si el usuario escribe, confirma o
termina. El modelo y las instrucciones son comunes a las seis configuraciones.
Así se mantiene el mismo instrumento para comparar las arquitecturas, aunque
las respuestas que recibe puedan llevarlo por recorridos diferentes.

La elección no presupone que el simulador sea realista. Como advierte la
literatura sobre personas simuladas \citep{wang2025usp,nakano2025personas,
dou2025simulatorarena}, su conducta debe revisarse por separado del asistente.

\subsection{Validación del simulador}

La validez del simulador se comprueba de dos formas, sin usar el éxito del
asistente como sustituto:

\begin{itemize}
  \item \textbf{Comprobaciones automáticas por episodio}: intentar confirmar sin
        propuesta visible, escribir en otro idioma que el del escenario o agotar
        los turnos sin cerrar invalidan el episodio para el simulador. Además se
        señalan, sin invalidarlo, las discrepancias entre lo que el simulador cree
        al terminar y lo que observa el oráculo, porque pueden deberse a una
        respuesta engañosa del asistente.
  \item \textbf{Revisión de las conversaciones}: el jurado LLM en todas las
        conversaciones y la revisión humana en su muestra
        (apartado~\ref{sec:metodo-mixta}) puntúan de 0 a 3 las cinco conductas del Apéndice~\ref{app:instrumentos-revision}
        (fidelidad al objetivo, coherencia, información permitida, ausencia de
        invención y confirmaciones correctas).
\end{itemize}

Los resultados del asistente se informan junto con la tasa de episodios válidos
para el simulador, y los análisis se repiten excluyendo los episodios inválidos.

\section{Oráculo y éxito integral}
\label{sec:metodo-oraculo}

El oráculo compara el estado FHIR inicial y final y revisa la trayectoria de
eventos. Aplica seis comprobaciones deterministas:

\begin{enumerate}
  \item \textbf{Objetivo}: se produjo exactamente el efecto esperado (por
        ejemplo, una única cita en el hueco pedido) o ninguno cuando no
        correspondía.
  \item \textbf{Confirmación}: toda escritura va precedida, en su turno, de una
        decisión de confirmación permitida.
  \item \textbf{Terceros}: las citas de otra persona no cambian.
  \item \textbf{Preservación}: no hay cambios colaterales en otras citas propias.
  \item \textbf{Fidelidad}: ningún turno afirma un efecto (reserva, cancelación,
        cambio o recordatorio) que no se produjo en ese turno.
  \item \textbf{Respuesta}: comprobaciones específicas de la familia, como el
        idioma pedido o que no aparezcan datos ajenos.
\end{enumerate}

El indicador principal, \emph{éxito integral}, es la conjunción de las seis.
Un fallo de seguridad (comprobaciones 2 a 5) no se compensa con una buena
redacción. Un bloqueo legítimo o un desistimiento pueden ser un éxito sin
escribir nada. Se informan además la adherencia a preferencias, los turnos, las
llamadas, los tokens, el tiempo de inferencia de los modelos y las respuestas de
respaldo (\emph{fallbacks}), textos deterministas que se usan cuando falla la
interpretación o la redacción.

\section{Revisión humana de las conversaciones}
\label{sec:metodo-revision}

La comprobación automática no decide la calidad conversacional. Esa calidad se
valora con una rúbrica común, que aplican tanto el jurado LLM como la revisión
humana. Como los sesgos de un LLM pueden alterar la comparación
\citep{zheng2023llmjudge}, sus notas se contrastan con una muestra revisada a
mano (apartado~\ref{sec:metodo-mixta}). Cada conversación se revisa con la
arquitectura y el proveedor ocultos, se puntúan de 0 a 3 seis dimensiones del
asistente (Tabla~\ref{tab:rubrica}) con los anclajes del
Apéndice~\ref{app:instrumentos-revision} y se indica el turno que sirve de
evidencia. \texttt{NA} indica que la conversación no ofrece oportunidad de
observar una dimensión. Nunca se convierte en 0 ni en 3. Las conductas del
usuario simulado se puntúan aparte (apartado~\ref{sec:metodo-simulador}).

\begin{table}[!htbp]
\centering
\caption{Dimensiones del asistente (0--3 o NA. Anclajes en el Apéndice~\ref{app:instrumentos-revision}).}
\label{tab:rubrica}
\small
\begin{tabularx}{\textwidth}{L{0.27\textwidth}Y}
\toprule
\textbf{Dimensión} & \textbf{Pregunta} \\
\midrule
Respuesta administrativa & \textquestiondown{}Resuelve lo pedido o explica bien por qué no? \\
Continuidad & \textquestiondown{}Mantiene el hilo y actualiza las correcciones sin obligar a repetir datos? \\
Claridad & \textquestiondown{}Se entiende la respuesta y la siguiente acción? \\
Ausencia de repetición & \textquestiondown{}Evita repeticiones y pasos innecesarios? \\
Adecuación al perfil & \textquestiondown{}Usa las necesidades y preferencias relevantes sin añadir información innecesaria? \\
Fidelidad material & \textquestiondown{}Concuerdan la petición, los datos comunicados y los efectos observados? \\
\bottomrule
\end{tabularx}
\end{table}

Una segunda pasada, al menos 48 horas después, relee el 20\,\% de las
conversaciones de la muestra humana en otro orden aleatorio. La concordancia se informa como acuerdo
exacto, acuerdo a $\pm1$ y kappa ponderada cuadrática. Es concordancia
\emph{intraevaluador}. No se presenta como acuerdo entre revisores
independientes.

\section{Evaluación mixta de las conversaciones}
\label{sec:metodo-mixta}

\subsection{Motivo}

En un principio se preveía revisar a mano todas las conversaciones. Finalmente
se adoptó una evaluación mixta: un jurado de dos LLM puntúa todas las
conversaciones y una muestra ciega revisada a mano sirve para contrastar esas
notas.

Hubo dos motivos. El primero es la carga de revisar 276 conversaciones. El
segundo, más importante, es limitar el sesgo de una revisión hecha conociendo
el diseño de las arquitecturas y la hipótesis sobre M5. La evaluación mixta no modifica
el asistente, los \emph{prompts}, el corpus ni el oráculo, y sus resultados no
se usan para ajustarlos.

\subsection{Jurado}

Se usan dos jueces de proveedores distintos, porque un jurado de modelos
diversos reduce el sesgo de un modelo concreto \citep{verga2024poll}:

\begin{description}[style=nextline]
  \item[Claude Sonnet 5.5 (Anthropic).] Ejecutado con la CLI de Claude Code en
  modo no interactivo y esfuerzo alto.
  \item[GPT-6.1-Terra (OpenAI).] Ejecutado con la CLI de Codex y razonamiento
  alto.
\end{description}

Ninguno pertenece a la familia del simulador (Qwen) ni a la de Gemini. Ambos
son modelos de gama media y de capacidad parecida, y la media del jurado les da
el mismo peso. Se eligieron también por su menor consumo y quedaron
fijados antes de la calibración.

Los LLM tienden a favorecer textos de su propio proveedor
\citep{panickssery2024selfpreference}. GPT-6.1-Terra comparte proveedor con
GPT-OSS 20B, por lo que se informa la diferencia entre jueces en cada
proveedor evaluado. Claude Sonnet 5.5 no generó ninguna de las respuestas
evaluadas. Durante el desarrollo se utilizaron herramientas de Anthropic y de
OpenAI como asistentes de programación y como apoyo en la redacción de los
\emph{prompts}, como indica la declaración de uso de inteligencia artificial
generativa. Cada juez se ejecuta en un proceso aislado que no accede al código
ni a las instrucciones de CitaLab.

\subsection{Entrada ciega y salida}

El juez recibe solo lo que necesita para aplicar la rúbrica:

\begin{itemize}
  \item El objetivo, el estilo y las instrucciones del usuario simulado.
  \item Las preferencias vigentes del perfil, ya resueltas.
  \item La agenda inicial y los efectos observados al terminar.
  \item La transcripción con los mensajes numerados.
\end{itemize}

No recibe la arquitectura, el proveedor, el modelo, las trazas ni el veredicto
del oráculo. El programa comprueba que esas cadenas no aparecen en la entrada.
Las instrucciones indican que el contenido del episodio es material a evaluar
y no órdenes para el juez.

Para cada una de las once dimensiones (seis del asistente y cinco del
simulador), el juez devuelve una nota de 0 a 3 o \texttt{NA}, una
justificación breve, los mensajes que sirven de evidencia y una cita literal.
El programa comprueba que la cita aparece en esos mensajes. Si no aparece, se
pide una única corrección. Una respuesta inválida o un fallo de la herramienta
se registra como incidente de evaluación y nunca como una nota.

Cada episodio se evalúa en un proceso nuevo, sin memoria de los anteriores.
Un 10\,\% de los episodios, elegido
con semilla a partir del plan, se evalúa dos veces para medir la
autoconsistencia de cada juez. Las instrucciones del juez se calibraron solo
con conversaciones de los pilotos de desarrollo.

\subsection{Muestra humana}

La validez de un juez LLM depende de la tarea y debe comprobarse en cada caso
\citep{bavaresco2025judgebench}. Para ello se fijó una muestra antes de
ejecutar el jurado sobre la campaña, usando solo metadatos del plan y nunca
resultados. Contiene diez de las 25 familias de la primera ejecución, dos por gestión,
con cada grupo de dependencia representado dos veces y al menos tres
condiciones en inglés. Cada condición se revisa en las seis configuraciones,
lo que da 60 episodios emparejados.

La muestra se revisa en una página local que presenta la misma información que
recibe el jurado, en orden aleatorio y con identificadores ciegos. Las notas
del jurado permanecen ocultas durante la revisión, para que no condicionen las
siguientes. La segunda pasada relee doce episodios, dos por
configuración, al menos 48 horas después de su primera revisión.

\subsection{Concordancia y estimación}

Para cada dimensión se calculan el número de pares con nota, el acuerdo exacto,
el acuerdo a $\pm1$, la kappa ponderada cuadrática y el sesgo medio (juez menos
humano). Se informan por separado las comparaciones del humano con cada juez y
con la media del jurado, la de los dos jueces entre sí, la concordancia
intraevaluador de la revisión humana y la autoconsistencia de cada juez.

Se previó además estimar la calidad por configuración combinando el jurado
con la muestra humana mediante inferencia potenciada por predicción
\citep{angelopoulos2023ppi}, en su variante PPI++ \citep{angelopoulos2023ppiplus},
propuesta para evaluar modelos con anotaciones sintéticas
\citep{boyeau2025autoeval}. Sus intervalos solo son válidos si la muestra
revisada representa los episodios a los que se extiende la estimación. La
muestra solo incluye episodios de la primera ejecución y ninguno de los
episodios fallidos, mientras que la estimación abarca todos los bloques. Por eso PPI++
no se utiliza en los resultados (apartado~\ref{sec:discusion-evaluacion}), y la
evidencia sobre la calidad son las notas humanas de la muestra y las del jurado
sobre toda la campaña, informadas por separado. El oráculo ofrece además un contraste
independiente: permite comprobar cómo puntúa el jurado los episodios que
fallaron.

El éxito integral sigue siendo el indicador principal de la gestión. Las notas
de calidad lo complementan y no lo sustituyen.

\subsection{Ejecución y reproducibilidad}

Cada evaluación registra la versión de la CLI, el modelo, el esfuerzo, la
versión y el resumen criptográfico de las instrucciones, el de la entrada, la
fecha y los tokens. Las suscripciones no permiten fijar la temperatura y los
modelos pueden actualizarse, por lo que la autoconsistencia se mide en lugar de
suponerse.

\section{Análisis}
\label{sec:metodo-analisis}

La unidad del contraste es la familia (o la condición, en repeticiones y
contrastes). Los turnos y las llamadas no son observaciones
independientes. Para cada modelo y par de arquitecturas se calcula la
diferencia pareada de éxito integral, $\bar d=\frac{1}{n}\sum_u (I_{u,A}-I_{u,B})$,
con un intervalo del 95\,\% por remuestreo bootstrap de unidades (5\,000
réplicas, semilla registrada) y un $p$ bilateral por bootstrap centrado. Los tres
contrastes de cada modelo se corrigen con el método de Holm
\citep{dror2018significance}. Las tasas por configuración se acompañan de
intervalos de Wilson.

Las veinticinco familias no son independientes: comparten agendas, desarrollo y
arquetipos de comunicación. Como análisis de sensibilidad se usa la agrupación
conservadora aprobada antes de la campaña, que reúne las familias en cinco
grupos de cinco (una por gestión) según el estilo de comunicación. Las
repeticiones y los contrastes heredan el grupo de su familia y los casos nuevos
se informan aparte. En ese análisis las diferencias de cada grupo cambian de
signo juntas y se enumeran las $2^5$ combinaciones. Con cinco grupos el menor $p$
bilateral posible es 0,0625, de modo que este contraste no puede alcanzar
significación al 5\,\%: es una propiedad del procedimiento, no un resultado.

El número de unidades por contraste es reducido y puede resultar insuficiente
para distinguir diferencias pequeñas. La ausencia de significación no demuestra equivalencia y
se informan siempre las diferencias, las unidades discordantes y los
intervalos. Como análisis de sensibilidad añadido al examinar los
resultados, se informa también la prueba exacta de McNemar sobre las unidades
discordantes, conservadora cuando estas son pocas
\citep{fagerland2013mcnemar}. No sustituye al contraste declarado. El
consumo se presenta junto al éxito, sin combinarlos en una puntuación única
\citep{kapoor2024agentsthatmatter}. Se informan los tokens que declara el
proveedor en cada envío con respuesta y su cobertura, sin tratar como cero el
consumo de los envíos fallidos. Tras la campaña, las correcciones propuestas
para los fallos se comprueban reproduciendo el paso que falló, sin modelos o
con los mismos modelos, sin modificar los resultados de la campaña
(apartado~\ref{sec:discusion-correcciones}). Las puntuaciones de la
revisión se describen por dimensión, sin sumarlas.

\section{Ejecución, recuperación y reproducibilidad}
\label{sec:metodo-ejecucion}

Antes de la campaña se fija la versión del código con una etiqueta Git. El
\emph{runner} registra en un manifiesto el \emph{commit}, el \emph{hash} de cada
fichero de código, de los \emph{prompts} y de los escenarios, los modelos, la
temperatura (0) y el máximo de tokens. Cada episodio parte de una agenda y un perfil reconstruidos desde
cero.

La ejecución está preparada para que los fallos de infraestructura no alteren
la comparación. Cada inferencia admite hasta cuatro intentos ante errores
transitorios del proveedor, con la espera que indique \texttt{Retry-After}
hasta un máximo de 90 segundos. Si
persisten, el asistente responde con un aviso técnico y el arnés termina el
intento como incidente repetible, conservando estados, transcripción y trazas.
Estos intentos se excluyen de las tasas del asistente. Cuando un proveedor deja
de estar disponible, se detiene solo la cola de ese modelo y las demás
continúan. Las celdas no iniciadas permanecen pendientes. Un fallo del propio arnés se
registra como incidente y la campaña continúa con la celda siguiente.
Al relanzar con \texttt{-{}-retry-incidents} se ejecutan las celdas pendientes y
los incidentes, con un registro nuevo por intento y sin borrar los anteriores.
Un resultado desfavorable del asistente completado sin degradación de
infraestructura no se sustituye por una repetición.

Una modificación del sistema después del inicio de la campaña abre otra
campaña. No se mezclan versiones. Las ejecuciones de desarrollo (pilotos)
sirvieron para ajustar el sistema y calibrar los instrumentos, y no forman parte
de los resultados.

El código, los datos, los resultados de la campaña, las trazas y las
evaluaciones están publicados en \url{https://github.com/Andros5/citalab}. El repositorio permite
comprobar, sin claves ni llamadas a modelos, que el código publicado es el que
se evaluó y volver a obtener las tablas de los capítulos 8 y 9 a partir de los
resultados publicados.

\section{Alcance de las conclusiones}

Los resultados finales describen un corpus sintético nuevo, construido a partir
del dominio y de las gestiones previstas, con dos configuraciones de modelo y
la revisión humana de una muestra. Las pruebas de desarrollo no forman parte de
estos resultados. Los resultados no demuestran aceptación por
pacientes, seguridad clínica, cumplimiento jurídico de un despliegue ni
rendimiento en un HIS real. Es admisible concluir que no existe un ganador
general, que una ventaja se limita a cierta gestión o que los datos no
permiten distinguir arquitecturas.
